\documentclass[11pt]{article}
\usepackage[utf8]{inputenc}
\usepackage[T1]{fontenc}
\usepackage{graphicx}
\usepackage{booktabs}
\usepackage{amsmath}
\usepackage{amssymb}
\usepackage[margin=1in]{geometry}
\usepackage{caption}
\usepackage{array}
\begin{document}
\title{Kelp-9 lozenges and flare-free survival in linen wheeze}
\author{Mara Pell and Ivo Quay}
\maketitle
\begin{abstract}
Adults with linen wheeze were followed at Amber Quay Infirmary and Saltglass Day Ward.
{{v:n_screened}} participants were screened, {{v:n_excluded}} were excluded before randomization, and {{v:n_randomized}} participants were randomized.
Kelp-9 was allocated to {{v:n_alloc_k}} adults and the blank lozenge was allocated to {{v:n_alloc_b}} adults.
Flare-free survival at {{v:m12}} months was {{v:s_k12}} (95\% CI {{v:lo_k12}} to {{v:hi_k12}}) percent with Kelp-9 and {{v:s_b12}} (95\% CI {{v:lo_b12}} to {{v:hi_b12}}) percent with the blank lozenge.
The difference was {{v:diff12}} (95\% CI {{v:lo_diff}} to {{v:hi_diff}}) percentage points.
\end{abstract}
\section{Introduction}
Linen wheeze is a recurring spell of tight breath and a dry whistle. A lozenge is sometimes started after the spell settles, in the hope of delaying the next spell \cite{pell2016,quay2018}.
Grouped follow-up, one product-limit step for each band, and a Greenwood interval on the probability scale are the pieces used to report that delay \cite{amber2019,harrow2017,moth2020}.
Printed percentages in this report are rounded half up \cite{brindle2021}.
The trial compared a Kelp-9 lozenge with a blank lozenge of the same size and color. Follow-up was planned through {{v:m12}} months at Amber Quay Infirmary and Saltglass Day Ward.
\section{Methods}
Adults were eligible when a linen wheeze spell had settled within the previous month and the adult could take a lozenge.
Randomization was one to one, from a shuffled list held at the ward desk.
The event was the return of a linen wheeze spell. Time was recorded in four bands, ending at months {{v:m3}}, {{v:m6}}, {{v:m9}}, and {{v:m12}}.
An adult who stopped attending during a band, and an adult who left after a band with no spell, was censored in that band.
The number at risk is the count still under observation at the opening of a band.
The number at risk for the next band equals that count minus the events in the band minus the adults censored in the band.
Each band is one step. Survival at the end of a band is the running product of (entered minus events) divided by entered.
The Greenwood variance is that survival, squared, times the sum of events divided by the product of entered and (entered minus events).
The standard error is the square root of the variance.
A 95\% interval on the probability scale is the estimate minus or plus {{v:z_mult}} standard errors.
The difference in {{v:m12}}-month survival uses the same form. Its variance is the sum of the two Greenwood variances.
Survival, standard errors, interval endpoints, and the difference are multiplied by 100 and then rounded half up to one decimal place.
Rounding is applied only at that last step, to the unrounded product and the unrounded standard error.
\section{Results}
{{v:n_screened}} participants were screened.
{{v:n_inelig}} adults did not meet the ward criteria and {{v:n_declined}} adults declined, so {{v:n_excluded}} adults were excluded before randomization.
{{v:n_randomized}} participants were randomized.
Kelp-9 was allocated to {{v:n_alloc_k}} adults and the blank lozenge was allocated to {{v:n_alloc_b}} adults.
{{v:n_miss_k}} adults in the Kelp-9 arm and {{v:n_miss_b}} adults in the blank arm missed the opening visit.
Follow-up opened for {{v:n_enter_k}} adults in the Kelp-9 arm and {{v:n_enter_b}} adults in the blank arm.
\falsifytable{T1}
In the Kelp-9 arm the opening counts were {{v:n_enter_k}}, {{v:n_k1}}, {{v:n_k2}}, and {{v:n_k3}}.
Events in those bands were {{v:d_k0}}, {{v:d_k1}}, {{v:d_k2}}, and {{v:d_k3}}, and the censored counts were {{v:c_k0}}, {{v:c_k1}}, {{v:c_k2}}, and {{v:c_k3}}.
In the blank arm the opening counts were {{v:n_enter_b}}, {{v:n_b1}}, {{v:n_b2}}, and {{v:n_b3}}.
Events were {{v:d_b0}}, {{v:d_b1}}, {{v:d_b2}}, and {{v:d_b3}}, and the censored counts were {{v:c_b0}}, {{v:c_b1}}, {{v:c_b2}}, and {{v:c_b3}}.
Flare-free survival with Kelp-9 was {{v:s_k3}}\%, {{v:s_k6}}\%, {{v:s_k9}}\%, and {{v:s_k12}}\% at months {{v:m3}}, {{v:m6}}, {{v:m9}}, and {{v:m12}}.
With the blank lozenge the percentages were {{v:s_b3}}\%, {{v:s_b6}}\%, {{v:s_b9}}\%, and {{v:s_b12}}\%.
At {{v:m12}} months, flare-free survival was {{v:s_k12}} (95\% CI {{v:lo_k12}} to {{v:hi_k12}}) percent with Kelp-9 and {{v:s_b12}} (95\% CI {{v:lo_b12}} to {{v:hi_b12}}) percent with the blank lozenge.
The Greenwood standard errors were {{v:se_k12}} percentage points for Kelp-9 and {{v:se_b12}} percentage points for the blank lozenge.
The difference was {{v:diff12}} (95\% CI {{v:lo_diff}} to {{v:hi_diff}}) percentage points.
The Greenwood standard error of that difference was {{v:se_diff}} percentage points.
After month {{v:m12}}, {{v:n_end_k}} adults in the Kelp-9 arm and {{v:n_end_b}} adults in the blank arm were still under observation and had not had a flare.
\falsifytable{T2}
\falsifyfigure{F1}
\section{Discussion}
{{v:ph_concl_claim}}

The blank arm recorded more flares in every band, and the interval for the difference lies above zero.

{{v:ph_hedge}}

{{v:ph_limit}}

Both wards used the same band edges and the same lozenge handling \cite{quay2018,amber2019}.
\falsifybibliography
\end{document}
